% Pursuing Stacks (cotes 134-2 to 134-8): archive cote 134-4, site page 96 (PDF page 97), typescript page 183.
% typescript. Draft compared: ps3.tex:4907-4967.
% AI-transcribed by Claude Opus 5.5 (claude-opus-5-5, adaptive thinking, effort high), Message Batches API, batch ps2 (msgbatch_01NVwr1Yuo8qJf5rykeGBXPK), 2026-11-12; not yet checked by eye.
\scanpage{134-4/96}
\mspage{183}

\uline{\xout{\unsure{W}}aspheric}, or \uline{$\underline{W}$-aspheric}, if the corresponding map in (Cat) is . It
\margin{|}turns out that \xout{this is} (\unsure{beaause|because} of the localization axiom on $\underline{W}$) this
is equivalent with f being ``universally a $\underline{W}$-equivalence'' i.e. f being
``universally in $\underline{W}_A$ '' , namely that for any base-change $\unsure{G'}\hand{\to}G$ in $A\hat{}$ ,
the corresponding map in $A\hat{}$
\[ f' : F\mathrm{x}_G G' \hand{\longrightarrow} G' \]
be in $\underline{W}_A$ i.e. be a $\underline{W}$-equivalence. \unsure{Of} course, when this condition is
satisfied, then for any base change, f' will be, not only a $\underline{W}$-equiva-%
lence, but even $\underline{W}$-aspheric - thus we can say that f is ``universally''
$\underline{W}$-aspheric'' - where ``universally'' refers to \uline{base change in} $A\hat{}$ . This of
course does not mean (and here one has to be slightly cautious) that
the corresponding map in (Cat) is universally $\underline{W}$-aspheric (which \xout{\unsure{would}}
refers to arbitrary \uline{base change in} (Cat)). But this \xout{uncoherence of}
apparent uncoherence is of no practical importance as far as terminology
goes, as the \xout{\unsure{notion of a}}\inserted[typed]{\xout{\unsure{word}} word} ``$\underline{W}$-aspheric map in $A\hat{}$'' \xout{and} is wholly \xout{sufficient}
adequate and sufficient for naming the notion, without any need to re-%
place it by the more complicated and ambiguous name ``universally $\underline{W}$-asphe-%
ric '', which therefore will never be used. We even could rule out the
formal uncoherence, by using the words $\underline{W}_A$-equivalence, $\underline{W}_A$-aspheric
maps \hand{`}(which \deleted{implies}\inserted{are even} universally $\underline{W}_A$-aspheric maps, without any ambiguity
any longer), as well as $\underline{W}_A$-aspheric objects - replacing throughout $\underline{W}$
by $\underline{W}_A$ . In practical terms, I think that when working consistently with
a single given $\underline{W}$ , we'll soon enough drop anyhow both $\underline{W}$ and $\underline{W}_A$ in the
terminology and notations !

A last point which deserves some\xout{s} caution, is that for general A ,
there is no implication between the \xout{\illegible{1}} two asphericity properties of
an object F of $A\hat{}$, namely of F being $\underline{W}$-aspheric (i.e. the object $A_{/F}$ of
(Cat) being $\underline{W}$-aspheric\xout{)}\hand{,} i.e. the map \xout{$A_{/F}$ \illegible{2} (Cat) being in $\underline{W}$}
\[ \hand{(*)} \qquad A_{/F} \hand{\longrightarrow} e \]
in (Cat) being in $\underline{W}$\hand{)}, and the property that F $\hand{\to}$ e be $\underline{W}$-aspheric ,
namely the map
\[ \hand{(**)} \qquad A_{/F} \hand{\longrightarrow} A_{/\hand{e}}\hand{= A} \]
in (Cat) being aspheric (which also means that the products Fxa for a
in A are $\underline{W}$-aspheric objects of $A\hat{}$, i.e. the categories $A_{/F\mathrm{x}a}$ are $\underline{W}$-asphe-%
ric, i.e. the maps
\[ A_{/F\mathrm{x}a} \hand{\longrightarrow} e \]
in (Cat) are\inserted{$\underline{W}$-} aspheric\hand{)}. A third related notion, weaker than the last one,

\margin{96}
